\subsection{DNS缓存污染攻击}
为了提高域名解析的性能，减轻权威服务器侧的网络流量负载，递归服务器引入了缓存机制，即递归解析服务器在接收到查询的域名相应后，可将其中的信息存储在本地一段时间，缓存时间由权威服务器指定。
在同一网络环境内其他用户所要查询的域名信息已被缓存，那么递归解析服务器可以直接返回缓存数据，无需再向远程的权威服务器重复发送查询请求。

缓存污染攻击的目的是为了篡改域名在缓存中的资源记录，以进一步实现网络钓鱼、传播恶意软件或收集用户敏感信息等。该攻击能实现的根源在于域名协议报文缺乏消息完整性验证的检查。
攻击者通过发送精心构造的DNS响应报文来欺骗递归解析服务器，使其将错误的解析结果存储在缓存中（缓存保存有时间限制）。当递归解析服务器被污染，缩用使用该递归解析服务器的互联网用户在解析该域名时都会被错误重定向。

上述攻击的关键在于，攻击者发送的伪造DNS响应报文能被递归解析服务器所接受，这在于以下字段核验通过：
\begin{itemize}
    \item IP地址匹配：这一条符合，因为响应的源IP地址为权威服务器地址，目的地址为递归解析服务器地址，均已知；
    \item 端口匹配：响应报文的源端口为DNS服务端口，属于常见端口为53，但目的端口为递归解析服务器发送请求时的随机端口，该值未知，需要遍历猜测。
    \item TXID匹配：响应TXID于请求TXID相等，该值无法提前获知，需要遍历猜测。
    \item 问题区域匹配：通过攻击者自行发起DNS解析服务，从而问题区内容有攻击者掌握，这一条符合。
\end{itemize}
以上所有字段核验通过后还要求攻击者伪造的响应比权威服务器的真实响应先到达递归解析服务器。上述16位源端口和16位TXID待遍历，意味着总的遍历次数为
$$
2^{16} \times 2^{16} = 2^{32}
$$
即需要在权威服务器返回真实响应前做完$2^{32}$次枚举，并将所有伪造响应都发送到递归解析服务器，DNS污染攻击才能成功。